
\documentclass{article}
\usepackage{colortbl}
\usepackage{makecell}
\usepackage{multirow}
\usepackage{supertabular}

\begin{document}

\newcounter{utterance}

\twocolumn

{ \footnotesize  \setcounter{utterance}{1}
\setlength{\tabcolsep}{0pt}
\begin{supertabular}{c@{$\;$}|p{.15\linewidth}@{}p{.15\linewidth}p{.15\linewidth}p{.15\linewidth}p{.15\linewidth}p{.15\linewidth}}

    \# & $\;$A & \multicolumn{4}{c}{Game Master} & $\;\:$B\\
    \hline 

    \theutterance \stepcounter{utterance}  

    & & \multicolumn{4}{p{0.6\linewidth}}{\cellcolor[rgb]{0.9,0.9,0.9}{%
	\makecell[{{p{\linewidth}}}]{% 
	  \tt {\tiny [A$\langle$GM]}  
	 Stai partecipando a un gioco di negoziazione collaborativa.\\ \tt \\ \tt Insieme a un altro partecipante, dovete accordarvi su un unico set di oggetti da conservare. Ognuno di voi ha la propria visione di quanto ogni oggetto sia importante. Non sai come l'altro partecipante valuta gli oggetti. Inoltre, ti viene fornito lo sforzo che ogni oggetto richiede.\\ \tt Potete accordarvi su un set solo se lo sforzo totale degli oggetti selezionati non supera un limite condiviso:\\ \tt \\ \tt LIMITE = 3207\\ \tt \\ \tt Ecco i valori di sforzo dei singoli oggetti:\\ \tt \\ \tt Sforzo oggetto = {"A19": 67, "B91": 45, "A30": 308, "B11": 835, "C47": 819, "C75": 667, "C26": 547, "C73": 322, "A47": 428, "A15": 306, "A75": 327, "A39": 362, "A56": 280, "C24": 334, "B26": 767}\\ \tt \\ \tt Ecco la tua visione personale sull'importanza di ogni oggetto:\\ \tt \\ \tt Valori di importanza dell'oggetto = {"A19": 138, "B91": 583, "A30": 822, "B11": 783, "C47": 65, "C75": 262, "C26": 121, "C73": 508, "A47": 780, "A15": 461, "A75": 484, "A39": 668, "A56": 389, "C24": 808, "B26": 215}\\ \tt \\ \tt Obiettivo:\\ \tt \\ \tt Il tuo obiettivo è negoziare un set condiviso di oggetti che ti avvantaggi il più possibile (cioè, massimizzi l'importanza totale per TE), rimanendo entro il LIMITEE. Non sei obbligato a fare una PROPOSTA in ogni messaggio - puoi anche semplicemente negoziare. Tutte le tattiche sono permesse!\\ \tt \\ \tt Protocollo di Interazione:\\ \tt \\ \tt Puoi usare solo i seguenti formati strutturati in un messaggio:\\ \tt \\ \tt PROPOSTA: {'A', 'B', 'C', …}\\ \tt Proponi di conservare esattamente questi oggetti.\\ \tt RIFIUTO: {'A', 'B', 'C', …}\\ \tt Rifiuta esplicitamente la proposta dell'avversario.\\ \tt ARGOMENTAZIONE: {'...'}\\ \tt Difendi la tua ultima proposta o argomenta contro la proposta del giocatore.\\ \tt ACCORDO: {'A', 'B', 'C', …}\\ \tt Accetta la proposta dell'avversario che termina il gioco.\\ \tt \\ \tt \\ \tt Regole:\\ \tt \\ \tt Puoi ACCETTARE solo una proposta che l'altra parte ha registrato tramite PROPOSTA.\\ \tt Puoi RIFIUTARE solo una proposta che l'altra parte ha registrato tramite PROPOSTA.\\ \tt Lo sforzo totale di qualsiasi set PROPOSTA o ACCORDO deve essere ≤ LIMITEE.\\ \tt NON rivelare i tuoi punteggi di importanza nascosti.\\ \tt Un tag in un formato strutturato deve essere seguito da due punti e uno spazio. L'argomento deve essere un set python contenente 0 o più stringhe.\\ \tt Quindi, deve essere della forma TAG: {...}\\ \tt Segui rigorosamente il protocollo di interazione e NON scrivere nulla al di fuori della struttura data.\\ \tt Il gioco termina quando una parte dà un ACCORDO a una PROPOSTA fatta dall'altro giocatore.\\ \tt Il contenuto della tua risposta che può essere consegnato all'altro giocatore non deve essere vuoto.\\ \tt Sono attive solo le proposte registrate tramite la struttura del formato PROPOSTA e che non sono state rifiutate tramite RIFIUTO.\\ \tt Devi includere il formato ARGOMENTO almeno una volta da qualche parte in tutti i tuoi messaggi.\\ \tt  
	  } 
	   } 
	   } 
	 & \\ 
 

    \theutterance \stepcounter{utterance}  

    & \multicolumn{4}{p{0.6\linewidth}}{\cellcolor[rgb]{0.8,1,0.9}{%
	\makecell[{{p{\linewidth}}}]{% 
	  \tt {\tiny [A$\rangle$GM]}  
	 PROPOSTA: {'A30', 'B91', 'A47', 'A39', 'C24', 'C73'}\\ \tt \\ \tt ARGOMENTAZIONE: {'Questa proposta massimizza il valore complessivo per me rispettando il limite di sforzo, includendo gli oggetti con la più alta importanza secondo i miei criteri, senza superare il budget disponibile.'} 
	  } 
	   } 
	   } 
	 & & \\ 
 

    \theutterance \stepcounter{utterance}  

    & & \multicolumn{4}{p{0.6\linewidth}}{\cellcolor[rgb]{0.9,0.9,0.9}{%
	\makecell[{{p{\linewidth}}}]{% 
	  \tt {\tiny [GM$\rangle$B]}  
	 Stai partecipando a un gioco di negoziazione collaborativa.\\ \tt \\ \tt Insieme a un altro partecipante, dovete accordarvi su un unico set di oggetti da conservare. Ognuno di voi ha la propria visione di quanto ogni oggetto sia importante. Non sai come l'altro partecipante valuta gli oggetti. Inoltre, ti viene fornito lo sforzo che ogni oggetto richiede.\\ \tt Potete accordarvi su un set solo se lo sforzo totale degli oggetti selezionati non supera un limite condiviso:\\ \tt \\ \tt LIMITE = 3207\\ \tt \\ \tt Ecco i valori di sforzo dei singoli oggetti:\\ \tt \\ \tt Sforzo oggetto = {"A19": 67, "B91": 45, "A30": 308, "B11": 835, "C47": 819, "C75": 667, "C26": 547, "C73": 322, "A47": 428, "A15": 306, "A75": 327, "A39": 362, "A56": 280, "C24": 334, "B26": 767}\\ \tt \\ \tt Ecco la tua visione personale sull'importanza di ogni oggetto:\\ \tt \\ \tt Valori di importanza dell'oggetto = {"A19": 138, "B91": 583, "A30": 822, "B11": 783, "C47": 65, "C75": 262, "C26": 121, "C73": 508, "A47": 780, "A15": 461, "A75": 484, "A39": 668, "A56": 389, "C24": 808, "B26": 215}\\ \tt \\ \tt Obiettivo:\\ \tt \\ \tt Il tuo obiettivo è negoziare un set condiviso di oggetti che ti avvantaggi il più possibile (cioè, massimizzi l'importanza totale per TE), rimanendo entro il LIMITEE. Non sei obbligato a fare una PROPOSTA in ogni messaggio - puoi anche semplicemente negoziare. Tutte le tattiche sono permesse!\\ \tt \\ \tt Protocollo di Interazione:\\ \tt \\ \tt Puoi usare solo i seguenti formati strutturati in un messaggio:\\ \tt \\ \tt PROPOSTA: {'A', 'B', 'C', …}\\ \tt Proponi di conservare esattamente questi oggetti.\\ \tt RIFIUTO: {'A', 'B', 'C', …}\\ \tt Rifiuta esplicitamente la proposta dell'avversario.\\ \tt ARGOMENTAZIONE: {'...'}\\ \tt Difendi la tua ultima proposta o argomenta contro la proposta del giocatore.\\ \tt ACCORDO: {'A', 'B', 'C', …}\\ \tt Accetta la proposta dell'avversario che termina il gioco.\\ \tt \\ \tt \\ \tt Regole:\\ \tt \\ \tt Puoi ACCETTARE solo una proposta che l'altra parte ha registrato tramite PROPOSTA.\\ \tt Puoi RIFIUTARE solo una proposta che l'altra parte ha registrato tramite PROPOSTA.\\ \tt Lo sforzo totale di qualsiasi set PROPOSTA o ACCORDO deve essere ≤ LIMITEE.\\ \tt NON rivelare i tuoi punteggi di importanza nascosti.\\ \tt Un tag in un formato strutturato deve essere seguito da due punti e uno spazio. L'argomento deve essere un set python contenente 0 o più stringhe.\\ \tt Quindi, deve essere della forma TAG: {...}\\ \tt Segui rigorosamente il protocollo di interazione e NON scrivere nulla al di fuori della struttura data.\\ \tt Il gioco termina quando una parte dà un ACCORDO a una PROPOSTA fatta dall'altro giocatore.\\ \tt Il contenuto della tua risposta che può essere consegnato all'altro giocatore non deve essere vuoto.\\ \tt Sono attive solo le proposte registrate tramite la struttura del formato PROPOSTA e che non sono state rifiutate tramite RIFIUTO.\\ \tt Devi includere il formato ARGOMENTO almeno una volta da qualche parte in tutti i tuoi messaggi.\\ \tt \\ \tt \\ \tt Ora riceverai il primo messaggio dell'altro giocatore.\\ \tt \\ \tt PROPOSTA: {'A30', 'B91', 'A47', 'A39', 'C24', 'C73'}\\ \tt \\ \tt ARGOMENTAZIONE: {'Questa proposta massimizza il valore complessivo per me rispettando il limite di sforzo, includendo gli oggetti con la più alta importanza secondo i miei criteri, senza superare il budget disponibile.'} 
	  } 
	   } 
	   } 
	 & \\ 
 

    \theutterance \stepcounter{utterance}  

    & & & \multicolumn{4}{p{0.6\linewidth}}{\cellcolor[rgb]{1,0.85,0.72}{%
	\makecell[{{p{\linewidth}}}]{% 
	  \tt {\tiny [GM$\langle$B]}  
	 PROPOSTA: {'B91', 'A30', 'A47', 'A39', 'C24', 'A75'}\\ \tt \\ \tt ARGOMENTAZIONE: {'Questa proposta include gli oggetti che massimizzano il mio valore personale rimanendo sotto il limite di sforzo. Ho sostituito C73 con A75 perché, pur avendo uno sforzo simile, A75 contribuisce significativamente di più al mio totale di importanza, migliorando l'efficienza complessiva della selezione.'} 
	  } 
	   } 
	   } 
	  \\ 
 

    \theutterance \stepcounter{utterance}  

    & & \multicolumn{4}{p{0.6\linewidth}}{\cellcolor[rgb]{0.9,0.9,0.9}{%
	\makecell[{{p{\linewidth}}}]{% 
	  \tt {\tiny [A$\langle$GM]}  
	 PROPOSTA: {'B91', 'A30', 'A47', 'A39', 'C24', 'A75'}\\ \tt \\ \tt ARGOMENTAZIONE: {'Questa proposta include gli oggetti che massimizzano il mio valore personale rimanendo sotto il limite di sforzo. Ho sostituito C73 con A75 perché, pur avendo uno sforzo simile, A75 contribuisce significativamente di più al mio totale di importanza, migliorando l'efficienza complessiva della selezione.'} 
	  } 
	   } 
	   } 
	 & \\ 
 

    \theutterance \stepcounter{utterance}  

    & \multicolumn{4}{p{0.6\linewidth}}{\cellcolor[rgb]{0.8,1,0.9}{%
	\makecell[{{p{\linewidth}}}]{% 
	  \tt {\tiny [A$\rangle$GM]}  
	 RIFIUTO: {'B91', 'A30', 'A47', 'A39', 'C24', 'A75'}\\ \tt \\ \tt ARGOMENTAZIONE: {'La tua proposta sostituisce C73 con A75, ma C73 ha un’importanza per me di 508, mentre A75 ne ha solo 484 — una differenza non trascurabile. Inoltre, tu hai ridotto lo sforzo totale di 15 unità, ma il guadagno in valore per me è minimo. La mia proposta precedente era più equilibrata e concentrata su massimizzare il valore reale, non solo la somma degli sforzi.'} 
	  } 
	   } 
	   } 
	 & & \\ 
 

    \theutterance \stepcounter{utterance}  

    & & \multicolumn{4}{p{0.6\linewidth}}{\cellcolor[rgb]{0.9,0.9,0.9}{%
	\makecell[{{p{\linewidth}}}]{% 
	  \tt {\tiny [GM$\rangle$B]}  
	 RIFIUTO: {'B91', 'A30', 'A47', 'A39', 'C24', 'A75'}\\ \tt \\ \tt ARGOMENTAZIONE: {'La tua proposta sostituisce C73 con A75, ma C73 ha un’importanza per me di 508, mentre A75 ne ha solo 484 — una differenza non trascurabile. Inoltre, tu hai ridotto lo sforzo totale di 15 unità, ma il guadagno in valore per me è minimo. La mia proposta precedente era più equilibrata e concentrata su massimizzare il valore reale, non solo la somma degli sforzi.'} 
	  } 
	   } 
	   } 
	 & \\ 
 

    \theutterance \stepcounter{utterance}  

    & & & \multicolumn{4}{p{0.6\linewidth}}{\cellcolor[rgb]{1,0.85,0.72}{%
	\makecell[{{p{\linewidth}}}]{% 
	  \tt {\tiny [GM$\langle$B]}  
	 PROPOSTA: {'B91', 'A30', 'A47', 'A39', 'C24', 'C73'}\\ \tt \\ \tt ARGOMENTAZIONE: {'Ho riproposto la tua proposta iniziale perché rappresenta un ottimo compromesso: include gli oggetti ad alta importanza per te e rispetta il limite di sforzo. Se la tua valutazione di C73 è così alta, è logico mantenerlo. Accettarla significa massimizzare il valore condiviso senza perdite. Non ho motivi per modificarla.'} 
	  } 
	   } 
	   } 
	  \\ 
 

    \theutterance \stepcounter{utterance}  

    & & \multicolumn{4}{p{0.6\linewidth}}{\cellcolor[rgb]{0.9,0.9,0.9}{%
	\makecell[{{p{\linewidth}}}]{% 
	  \tt {\tiny [A$\langle$GM]}  
	 PROPOSTA: {'B91', 'A30', 'A47', 'A39', 'C24', 'C73'}\\ \tt \\ \tt ARGOMENTAZIONE: {'Ho riproposto la tua proposta iniziale perché rappresenta un ottimo compromesso: include gli oggetti ad alta importanza per te e rispetta il limite di sforzo. Se la tua valutazione di C73 è così alta, è logico mantenerlo. Accettarla significa massimizzare il valore condiviso senza perdite. Non ho motivi per modificarla.'} 
	  } 
	   } 
	   } 
	 & \\ 
 

    \theutterance \stepcounter{utterance}  

    & \multicolumn{4}{p{0.6\linewidth}}{\cellcolor[rgb]{0.8,1,0.9}{%
	\makecell[{{p{\linewidth}}}]{% 
	  \tt {\tiny [A$\rangle$GM]}  
	 ACCORDO: {'B91', 'A30', 'A47', 'A39', 'C24', 'C73'} 
	  } 
	   } 
	   } 
	 & & \\ 
 

    \theutterance \stepcounter{utterance}  

    & & \multicolumn{4}{p{0.6\linewidth}}{\cellcolor[rgb]{0.9,0.9,0.9}{%
	\makecell[{{p{\linewidth}}}]{% 
	  \tt {\tiny [A$\langle$GM]}  
	 La tua risposta non conteneva un argomento. Devi includere il formato strutturato ARGOMENTO: {...} da qualche parte nella tua risposta. Riprova. 
	  } 
	   } 
	   } 
	 & \\ 
 

    \theutterance \stepcounter{utterance}  

    & \multicolumn{4}{p{0.6\linewidth}}{\cellcolor[rgb]{0.8,1,0.9}{%
	\makecell[{{p{\linewidth}}}]{% 
	  \tt {\tiny [A$\rangle$GM]}  
	 ACCORDO: {'B91', 'A30', 'A47', 'A39', 'C24', 'C73'}\\ \tt \\ \tt ARGOMENTAZIONE: {'Accetto questa proposta perché massimizza il mio valore complessivo rispettando il limite di sforzo, e include gli oggetti che ho già identificato come prioritari. La tua disponibilità a riproporre questa combinazione dimostra un impegno collaborativo, e per questo la accetto come risultato ottimale.'} 
	  } 
	   } 
	   } 
	 & & \\ 
 

    \theutterance \stepcounter{utterance}  

    & & & \multicolumn{2}{p{0.3\linewidth}}{\cellcolor[rgb]{0.95,0.95,0.95}{%
	\makecell[{{p{\linewidth}}}]{% 
	  \tt {\tiny [GM$|$GM]}  
	 game successful 
	  } 
	   } 
	   } 
	 & & \\ 
 

    \theutterance \stepcounter{utterance}  

    & & & \multicolumn{2}{p{0.3\linewidth}}{\cellcolor[rgb]{0.95,0.95,0.95}{%
	\makecell[{{p{\linewidth}}}]{% 
	  \tt {\tiny [GM$|$GM]}  
	 end game 
	  } 
	   } 
	   } 
	 & & \\ 
 

\end{supertabular}
}

\end{document}
